\section{Evaluation Methodology}
\label{sec:evaluation-methodology}

The proposed SysSecOps process was evaluated through a controlled prototype case study
comprising three complementary investigations. The first investigation examined the
operational feasibility and execution reliability of the complete pipeline. The second
examined the applicability of the process to an existing on-premises virtual machine.
The third examined the performance of the machine-learning component and the correctness
of the unified risk-scoring mechanism. This design enabled the evaluation to address the
operation of the prototype, its applicability within a hybrid environment, and its
ability to transform heterogeneous security evidence into a reproducible deployment
decision.

The evaluation focused on observable system behaviour rather than on the capabilities of
any individual implementation tool. Evidence was collected from pipeline execution
records, synthesized infrastructure artifacts, scanner outputs, risk reports, deployment
records, target-state observations, and classifier predictions. The same scoring rules
and decision thresholds defined in Chapter~\ref{Chapter3} were applied throughout the
evaluation.

\subsection{Research Questions}
\label{subsec:evaluation-rqs}

The evaluation was guided by the following research questions:

\begin{description}
    \item[\textbf{RQ1.}] \textbf{Decision conformance and enforcement correctness:}
    When infrastructure of a known risk band is submitted, does the prototype
    reach the decision that band warrants, and does that decision actually
    govern what happens next --- such that a \emph{reject} never deploys, a
    \emph{review} never deploys without a recorded approval, and every run
    leaves a machine-readable account of itself, including when an optional
    assessment component or credential is unavailable?

    \item[\textbf{RQ2.}] \textbf{Cross-boundary enforcement equivalence:} When the
    same class of weakness is expressed once as public-cloud infrastructure code
    and once as on-premises configuration, does the shared risk engine
    \emph{detect} it on both sides of the hybrid boundary, and does the
    on-premises path attain the same operational checkpoints --- reachability,
    validation, gating, dry run, apply, post-apply verification, and repeated-run
    idempotency --- as the cloud path?

    \item[\textbf{RQ3.}] \textbf{Unified risk-evaluation capability:} To what extent
    can the proposed risk-scoring mechanism integrate findings from multiple security
    tools and a logistic-regression code-risk estimate to distinguish insecure code and
    produce consistent, interpretable deployment decisions?
\end{description}

Two aspects of this formulation need to be reviewed. First, RQ1 deliberately does
not ask whether runs \emph{complete}. A gate that never blocks anything
completes every run, so a completion rate cannot by itself distinguish a working
security control from an inert one. Conformance to a declared expectation, and
enforcement of the resulting decision, are the properties that carry evidential
weight; completion is reported only as a supporting measure. Second, RQ2 is
stated as a question about detection rather than about decision equality. An
earlier formulation asked only whether the process \emph{could be applied} to an
on-premises node, which the prototype's existence already answers. A later one
asked whether both sides reach the same decision band --- but the two sides are
not the same artifact, and demanding identical bands would assert an equivalence
between a firewall rule and a security group that does not hold, while being
satisfiable by simply retuning weights. What can fail, and what matters
operationally, is whether a weakness visible on one side is invisible on the
other. Section~\ref{subsec:evaluation-results-parity} reports a case where it
was.


Table~\ref{tab:rq-evaluation-map} presents the relationship between the research
questions, the corresponding evaluation methods, and the principal measurements.

\begin{table}[ht]
\centering
\small
\setlength{\tabcolsep}{4pt}
\caption{Relationship between the research questions and evaluation measures.}
\label{tab:rq-evaluation-map}
\begin{tabular}{|C{1.2cm}|L{4.3cm}|L{4.5cm}|L{4.1cm}|}
\hline
\textbf{RQ} & \textbf{Evaluation method} & \textbf{Collected evidence} &
\textbf{Measures} \\
\hline
RQ1 & Replicated execution of pre-registered scenarios spanning three risk bands
on both branches, plus scanner-removal scenarios & Gate reports, per-stage timing
records, enforcement outcomes, and campaign run records & Decision conformance,
confusion matrix, enforcement-invariant satisfaction, degradation behaviour, and
gate overhead \\
\hline
RQ2 & Paired cloud/on-premises fixtures expressing the same weakness class, and
checkpoint attainment on a Linux node reached over SSH & Per-branch finding
categories, paired gate scores, connectivity results, dry-run and
apply records, verification results, and idempotency counters & Detection
coverage per weakness class, checkpoint attainment rate, and failure-class
distribution \\
\hline
RQ3 & Held-out classifier evaluation and controlled risk-decision cases & True labels,
predicted probabilities, confusion matrix, component scores, deduplicated findings, and
gate decisions & Accuracy, precision, recall, F1-score, ROC--AUC, score correctness, and
decision consistency \\
\hline
\end{tabular}
\end{table}

\subsection{Evaluation Design}
\label{subsec:evaluation-design}

The unit of analysis for RQ1 was one complete pipeline run. A run began when an
infrastructure request or existing project was submitted to the prototype and ended when
the process produced a persisted terminal decision. Where the decision permitted
deployment, the run additionally included deployment and post-deployment status
collection. For RQ2, the unit of analysis was one pipeline-controlled Ansible operation
against an existing virtual machine. For RQ3, an individual Python file constituted the
unit of analysis for classifier evaluation, while one controlled combination of risk
inputs constituted the unit of analysis for gate-decision evaluation.

The evaluation varied three principal conditions: the risk characteristics of the
submitted infrastructure, the deployment target, and the availability of optional
assessment components. The observed dependent variables included completion status,
stage outcome, gate decision, deployment outcome, resulting target state, predicted
code-risk probability, and agreement between independently calculated and reported risk
scores.

A security rejection was not classified as a pipeline failure. A run that identified an
unacceptable change, produced a valid \emph{reject} decision, and prevented deployment
represented correct execution of the security process. An unexpected failure was defined
as an unhandled exception, premature pipeline termination, absent or malformed decision
report, incorrect transition between stages, or deployment that violated the gate and
approval policy.

\subsection{Pre-Registered Evaluation Campaign}
\label{subsec:evaluation-campaign}

The measurements reported in Section~\ref{sec:evaluation-results} were not
gathered by reading back the operational logs the prototype happened to
accumulate during development. Such logs are a convenience sample: they
over-represent whatever was being debugged at the time, they contain runs of
incomplete projects, and the analyst chooses after the fact which of them to
count. To avoid this, the evaluation was restructured around a campaign whose
scenarios, expected outcomes, and replicate counts were declared in a
version-controlled specification \emph{before} any of the reported runs were
executed, and whose results were then computed from the persisted artifacts by a
single analysis program.

\subsubsection{Calibrated Fixtures}

Thirty infrastructure fixtures were constructed. Six are \emph{band fixtures}: for each branch
--- AWS CDK for the cloud path and Ansible for the on-premises path --- one
fixture was authored to sit in each of the \emph{pass}, \emph{review}, and
\emph{reject} bands. The remaining twenty-four form twelve \emph{matched pairs},
each pair expressing the same thing twice, once in CDK and once in Ansible.
Eleven pairs express one catalogued weakness class each: unrestricted network
ingress, a plaintext credential, over-privileged access, unencrypted storage,
insufficient logging, missing authentication, incorrect permissions, externally
accessible files, unbounded resources, cleartext transmission, and execution
with unnecessary privileges. The twelfth pair is the \emph{specificity control},
a compliant stack and a compliant playbook that express no weakness at all and
against which every detection matcher must remain silent.

Each fixture carries a declared band in \texttt{expectations.yaml} together with
the finding breakdown that produces it. A calibration utility re-scores every
fixture and fails if any has drifted outside its declared band, so a scanner
upgrade that silently changes the arithmetic is caught rather than absorbed. An
important methodological constraint applies to the matched pairs: their declared
bands are \emph{descriptive}, recorded from measurement, and were not adjusted
to bring the two branches into line. Tuning the fixtures or the weights until
both sides returned the same band would have manufactured a result while leaving
the question RQ2 untouched. For the same reason the bands recorded for the
coverage pairs are those measured in the baseline arm and were deliberately left
unchanged when the detection rules were added; the remediated arm consequently
reports the affected fixtures as non-conformant against them. That divergence is
the effect being measured, and updating the bands to match would erase it.

\subsubsection{Scenario Matrix}

The specification in \texttt{evaluation/scenarios.yaml} declares 48 scenarios
totalling 154 offline runs, organised in six groups. Eight \emph{conformance}
scenarios, replicated five times each, submit the six band fixtures and
additionally exercise the \emph{review} band twice per branch --- once withheld
and once with an operator approval supplied --- so that the approval gate is
observed in both states. Ten \emph{degradation} scenarios, replicated three
times each, remove one scanner from an otherwise identical run; each declares
the decision that the remaining evidence should support, derived arithmetically
from the calibrated breakdown, so the scenario tests a prediction rather than
merely recording whatever occurs. Eight \emph{parity} scenarios, replicated
three times each, submit the four original pairs, and four \emph{parity
baseline} scenarios repeat the on-premises half of each pair with the semantic
configuration rules disabled; these are retained as a regression check that the
rules are the component responsible for the on-premises findings, and are not
reported separately. Sixteen \emph{coverage} scenarios, replicated three times
each, submit the seven additional matched pairs derived from the external
weakness catalogue together with the two specificity controls; these and the
eight parity scenarios form the coverage arm of 24 scenarios and 72 runs, which
is executed twice --- once at the frozen gate commit and once after the
detection rules were added --- yielding the baseline and remediated
measurements. The conformance, degradation and parity groups together form the
106-run batch on which the RQ1 results are reported; the eight parity scenarios
belong to both that batch and the coverage arm, which is why the two totals
overlap. Two further scenarios,
replicated three times each, exercise real deployment against a live node; they
are skipped unless the \texttt{target-host} capability is explicitly granted,
which keeps the offline campaign self-contained and reproducible without any
external dependency. They are reported separately as the live-target arm, run
against a disposable Multipass virtual machine provisioned by
\texttt{scripts/setup\_multipass\_target.sh}.

\subsubsection{Controlling Non-Determinism}

Two components make a run's score depend on the day it executes: the cost
estimator prices against a live feed, and the cloud configuration check reports
live account state. During pilot execution this was not hypothetical --- a
\emph{reject} fixture scored 119 against a calibrated 114 because live pricing
contributed an additional cost band. The registered campaign therefore disables
both components by default and records that fact in each run's manifest, so that
a reported score is reproducible from the repository alone. A separate,
explicitly non-deterministic arm can be requested when the interaction with live
pricing is itself of interest.

\subsubsection{Recorded Artifacts and Analysis}

Every run writes a self-describing report containing the run identifier,
wall-clock duration, terminal status, the scenario and campaign it belongs to,
the conditions in force --- deployment enabled, approval supplied, scanners
disabled, and the numeric thresholds --- the provenance of the environment, and
a per-stage record giving each stage's outcome, duration, and, on failure, a
classified failure reason. Provenance capture includes the repository commit,
whether the working tree was dirty, the interpreter version, the host, and the
resolved version of every external tool, so that a table can be traced to the
exact software state that produced it. The campaign runner appends one row per
run to a line-delimited record and copies each report into the campaign
directory.

All reported statistics are computed from those persisted records alone; the analysis program never re-executes the
pipeline. Proportions are reported as a count over the number of runs to which
the measure applied, rather than as a percentage alone. The denominator is
carried alongside every rate because several of them are small --- some
invariants and checkpoints are exercised by ten runs or fewer --- and a rate of
$1.00$ over ten runs supports a much weaker claim than the same rate over 106.
No confidence intervals are reported: the campaign's replicate counts were
chosen to cover the declared scenarios rather than to estimate failure
frequencies, and an interval computed over three or ten runs would suggest a
precision the design does not provide. The run counts are stated instead, so
that the strength of each result is visible directly. Durations are reported as
medians with inter-quartile ranges rather than means, because pipeline stage
times are right-skewed by occasional tool start-up costs.

\subsection{End-to-End Pipeline Evaluation}
\label{subsec:evaluation-rq1}

RQ1 was investigated by executing the prototype across representative scenarios that
exercised the principal paths through the operational workflow. Each run was observed at
the following checkpoints:

\begin{enumerate}
    \item acceptance or generation of the infrastructure source;
    \item synthesis of the public-cloud template or syntax validation of the
    private-side playbook;
    \item execution of the configured assessment components;
    \item normalisation and deduplication of the collected findings;
    \item calculation and persistence of the component scores, total score, and gate
    decision;
    \item enforcement of the deployment and approval policy; and
    \item publication of the decision and deployed-target state to the monitoring and
    audit layer.
\end{enumerate}

Four scenario classes were included. A low-risk change exercised the \emph{pass} path,
a moderate-risk change exercised the manual \emph{review} path, and a deliberately
insecure change exercised the \emph{reject} path. A degraded-execution scenario was also
included by making an optional assessment dependency or credential unavailable. This
scenario examined whether the corresponding component returned an explicit unavailable
or skipped status without causing an unhandled termination of the complete pipeline.
The degraded scenario was analysed separately because successful continuation with a
missing assessment source provides less security evidence than a fully assessed run.

For $n$ pipeline runs, the end-to-end completion rate was calculated as

\begin{equation}
\label{eq:pipeline-completion-rate}
\mathrm{Completion\ Rate} =
\frac{n_{\mathrm{terminal}}}{n} \times 100\%,
\end{equation}

where $n_{\mathrm{terminal}}$ denotes the number of runs that reached a valid and
persisted terminal decision without an unexpected failure. This rate is reported
as a supporting measure only. The primary measure for RQ1 is \emph{decision
conformance}, the proportion of runs whose observed decision equals the decision
declared for that scenario in advance:

\begin{equation}
\label{eq:decision-conformance}
\mathrm{Conformance} =
\frac{|\{r : d_{\mathrm{obs}}(r) = d_{\mathrm{exp}}(r)\}|}{n} \times 100\%,
\end{equation}

where $d_{\mathrm{obs}}(r)$ and $d_{\mathrm{exp}}(r)$ denote the observed and
pre-declared decision for run $r$. Conformance is reported both in aggregate and
as a confusion matrix over the three decision bands,
so that a systematic bias towards permissiveness or towards over-blocking would
be visible rather than averaged away.

Stage-level completion was reported in addition to the aggregate rate to identify
whether failures occurred during generation, validation, assessment, scoring,
deployment, or monitoring. Enforcement was then evaluated against a set of
invariants that must hold for every run regardless of its decision: a
\emph{reject} decision is never followed by a deployment; a \emph{review}
decision is followed by a deployment only when an approval record exists; every
run persists a decision report; every non-passing decision persists a
corresponding rejection or approval record; and every report carries the
provenance needed to reproduce it. Each invariant is reported as a satisfaction
rate over the runs to which it applies, so a single violation is visible rather
than diluted by the aggregate.

Finally, gate overhead was measured, since a control that is correct but
prohibitively slow will not survive contact with a real delivery process. The
duration of the gate stage and its share of total run wall-clock time were
recorded per run and reported as medians with inter-quartile ranges.

\subsection{Cross-Boundary Equivalence Evaluation}
\label{subsec:evaluation-rq2}

RQ2 examined whether the two sides of the hybrid boundary are governed alike.
This has two parts, which are measured separately: whether the on-premises path
can attain the same operational checkpoints as the cloud path, and whether the
shared risk engine detects the same class of weakness when it is expressed on
either side.

The on-premises target was a Linux virtual machine represented in an Ansible
inventory and configured through an Ansible playbook. For the live-target arm it
was a disposable Ubuntu~22.04 instance provisioned through Multipass and reached
over SSH with a purpose-generated key pair injected at first boot; the hybrid
case study of Section~\ref{sec:implementation} instead reaches an equivalent
private-side host over a mesh-VPN connection between the two domains. The
distinction is one of transport only --- in both cases the orchestrator sees an
inventory entry it must reach, authenticate to, and converge --- and using a
disposable instance keeps the measurement reproducible by anyone with the
provisioning script. The same risk engine, scoring rules, and decision
thresholds applied to public-cloud changes governed this private-side target.

The inventory is supplied as a file declaring the host in the
\texttt{appservers} group rather than as a bare hostname on the command line.
This is a correctness requirement, not a preference: the fixture playbooks target
that group, and \texttt{ansible-playbook} exits successfully with an empty
\texttt{PLAY RECAP} when no host matches a play, so a group mismatch would record
a deployment that applied nothing as a successful apply.

\subsubsection{Checkpoint Attainment}

Seven operational checkpoints were defined, corresponding to the stages a change
must pass to move from submission to a verified deployed state: target
reachability, syntax validation, security gating, dry run, apply, post-apply
verification, and repeated-run idempotency. For each run, a checkpoint was
recorded as attained, failed, or not attempted. Attainment rates were computed
over attempted checkpoints only. This distinction matters: when the gate returns
\emph{reject}, the dry-run checkpoint is never attempted, and counting that as a
failure would penalise the system for behaving correctly.

The evaluation procedure comprised six stages. First, network reachability and
authenticated Ansible access to the target were verified. Second, the inventory and
playbook were submitted to syntax validation and static security assessment. Third, the
resulting findings were normalised and processed by the same risk engine and decision
thresholds used for public-cloud changes. Fourth, playbook execution was permitted only
when the gate decision and approval state satisfied the deployment policy. Fifth, the
requested configuration state was verified directly on the virtual machine. Finally,
the target identity, gate decision, execution result, and post-execution state were
recorded for operational traceability.

Successful integration required all of the following conditions: authenticated access
to the existing target, successful analysis of the playbook, production of a valid gate
decision, enforcement of that decision before execution, successful application of an
allowed configuration, and verification of the requested target state. Where the
playbook described an idempotent operation, a subsequent execution was used to determine
whether the desired state remained stable without unintended configuration changes.

Observed failures were classified as connectivity, authentication, validation,
assessment, gate-enforcement, configuration-execution, or target-verification failures.
This classification separated failures caused by the external infrastructure from those
caused by the implementation of the proposed process.

\subsubsection{Detection Coverage}

Checkpoint attainment establishes that the on-premises path \emph{runs} the same
workflow. It does not establish that the workflow \emph{sees} the same things.
Testing the latter requires a set of weakness classes to test over, and the
choice of that set is the point at which such an evaluation is most easily
compromised. A set chosen by the author cannot distinguish a gate that covers
the space of deployment-time weaknesses from a gate whose author chose the
weaknesses it already covers. The class list was therefore derived from an
external catalogue rather than written for the purpose.

\paragraph{Derivation of the class list.}
The catalogue is MITRE's Common Weakness Enumeration view CWE-1008,
\emph{Architectural Concepts}, which organises weaknesses by the security
principle they violate rather than by language or platform \cite{MITRE_CWE1008}.
That organisation is what makes it usable here: a view indexed by principle
yields classes that can be expressed on either side of the hybrid boundary,
whereas a view indexed by technology would not. Version 4.20 of the view was
retrieved as the published XML export, and its archive digest, version, and
release date are recorded so
that the derivation can be repeated against the same source.

The view contains 12 categories spanning 223 member weaknesses. Three filters
were applied mechanically.
The first retains only entries that are members of a CWE-1008 category. The
second retains only entries whose abstraction is \emph{Base} or \emph{Class},
discarding variants too specific to have an expression in both technologies and
pillars too abstract to be tested at all. The third retains only entries whose
declared Modes of Introduction include at least one of \emph{Operation},
\emph{System Configuration}, \emph{Installation}, or \emph{Bundling} --- that
is, weaknesses that a deployment artifact can introduce, as distinct from those
introduced only while writing application code. A gate that inspects
infrastructure definitions cannot be asked to see the latter, and including them
would understate coverage by counting weaknesses outside the mechanism's remit.

The filters reduce 223 members to 38 candidates spanning seven categories; five
categories are eliminated entirely, and the reason for each elimination is
recorded rather than left implicit, because a category with no deployment-time
expression bounds what any pre-deployment gate can be asked to see. From the 38
candidates, 11 classes were retained as those expressible as a matched pair in
both a CDK application and an Ansible playbook; the 27 candidates not retained
are listed individually with the reason for exclusion. The resulting catalogue
is a pre-registration: it fixes the classes, their CWE identifiers, and the
matcher that decides whether a report names each class, before any of them was
measured.

\paragraph{Fixtures and the specificity control.}
Each of the 11 classes was authored twice --- once as CDK infrastructure code
and once as an Ansible playbook --- forming a matched pair in which the intended
weakness is held constant and only the boundary varies. Each fixture carries one
weakness and is otherwise deliberately unremarkable, so that a difference in the
gate's report is attributable to the weakness rather than to incidental
differences between the two artifacts.

Crediting detection from the weakness fixture alone would reward a matcher that
responds to everything. Each branch therefore also submits a \emph{control}
fixture carrying none of the catalogued weaknesses, and a class is credited only
if its matcher fires on the weakness fixture \emph{and} stays silent on the
control for the same branch. Where the control was never examined, the class is
recorded as not measured rather than as detected, so that an absent control
cannot be mistaken for a silent one.

The controls were not the ones originally intended. The RQ1 \emph{pass} fixtures
were used first, and the cloud one proved unusable: it provisions a bucket that
receives S3 access logs, and such a bucket cannot log to itself and must carry a
log-delivery ACL, so the report named both access logging and a public ACL no
matter how the fixture was written. Two classes therefore had a control that
could not discriminate. Purpose-built controls were substituted --- on the cloud
side a key, an encrypted log group, and an encrypted queue with a dead-letter
queue; on the on-premises side a playbook that installs a service under a
dedicated account with stated limits --- and both were verified to produce no
findings at all before the measurement was taken. The original measurement is
retained rather than discarded.

\paragraph{Two arms.}
The gate was frozen at a recorded commit before the first fixture was authored,
and the \emph{baseline} arm was measured at that commit. Detection rules written
in response to the baseline result were then implemented, and the
\emph{remediated} arm was measured over the identical scenario set. Reporting both,
with the frozen commit recorded, states plainly which capability preceded the
measurement and which followed it.

Two constraints keep the second arm from collapsing into fitting. First, the
matchers in the catalogue are frozen at their pre-registered form even where
they are known to under-credit, so a class is earned by emitting the
pre-registered category rather than by widening the pattern that recognises it
after seeing what the pattern missed. Second, each new rule is accompanied by
unit tests that express the same weakness through a different module or resource
type from the one the fixture uses, and by a test asserting that the rule stays
silent on the secure form of the same construct. A rule that passed only against
its own fixture would satisfy neither.

\paragraph{What is measured, and what is not.}
The measured quantity is coverage, not agreement. For each pair and each branch,
the persisted gate report is examined for a finding whose category names the
weakness the fixture encodes; a generic style or hygiene finding that happens to
touch the same artifact is not credited, since a report that names the wrong
problem gives an operator no signal that the real one exists. Alongside the
per-class result, the classes are cross-tabulated by which branches saw them, so
that a gap is located on a specific side of the boundary rather than averaged
away.

Decision agreement was deliberately \emph{not} adopted as the measure. The two
sides of a pair are not the same artifact --- a host firewall rule and a
security group differ in blast radius, in reversibility, and in how many
distinct controls they can violate --- so requiring identical bands would assert
an equivalence that does not hold, and could be satisfied by adjusting weights
rather than by improving detection. Where a branch fails to detect a weakness,
the finding breakdown of both sides is reported so that the cause can be
attributed to a specific absence of capability rather than left as an
unexplained numeric difference.

\subsection{Risk-Scoring Evaluation}
\label{subsec:evaluation-rq3}

RQ3 was examined through two evaluations. The first measured the ability of the
logistic-regression component to discriminate between labelled insecure and reference
secure Python files. The second examined whether the complete gate correctly aggregated
the machine-learning result with the remaining risk components and produced the decision
specified by the scoring policy.

\subsubsection{Dataset Construction}
\label{subsubsec:evaluation-dataset}

The evaluation corpus contained a positive class and a negative reference class. The
positive class was obtained from the \textsc{SecurityEval} dataset, which contains
Python vulnerability examples associated with Common Weakness Enumeration categories
and was developed for evaluating the security of machine-learning-based code-generation
techniques \cite{Siddiq2022SecurityEval}. Each selected vulnerability example was stored
as an individual Python file and assigned label $1$.

The negative class was sampled from the source trees of maintained open-source Python
projects, including \texttt{click}, \texttt{rich}, and \texttt{typer}. Each selected
file was assigned label $0$. These files constituted a \emph{reference secure} class
rather than a formally verified secure class. Their use in maintained software projects
provided practical negative examples, but did not establish the absence of every
possible vulnerability.

To improve comparability between the two classes, package initialisers and test files
were excluded. Files containing fewer than 20 or more than 700 non-blank,
non-comment lines were also excluded. These criteria reduced the influence of trivial
snippets, tests, and unusually large files on the learned model. The number of samples,
class distribution, and originating project or dataset were retained as dataset
metadata. The final corpus was balanced as far as practical to limit majority-class
bias.

\subsubsection{Feature Extraction and Model Training}
\label{subsubsec:evaluation-model-training}

Each file was analysed independently using Bandit and Semgrep. The findings were reduced
to an eleven-dimensional feature vector comprising Bandit severity counts, Bandit
confidence counts, Semgrep severity counts, and the total number of findings returned by
each analyser~\ref{subsubsec:rs-features}.

The corpus was divided into training and test partitions using a stratified $80/20$
split. Stratification preserved the class distribution in both partitions. A feature
scaler was fitted exclusively on the training partition and was then applied unchanged
to the held-out test partition. The logistic-regression classifier was trained with
balanced class weights, and its output for each file was the estimated probability of
membership in the insecure class.

\subsubsection{Classifier Evaluation Measures}
\label{subsubsec:evaluation-classifier-measures}

Classifier performance was assessed using the confusion matrix, accuracy, precision,
recall, F1-score, and the area under the receiver operating characteristic curve
(ROC--AUC). Let $TP$, $TN$, $FP$, and $FN$ denote the numbers of true positives, true
negatives, false positives, and false negatives, respectively. The principal measures
were calculated as

\begin{align}
\mathrm{Accuracy} &= \frac{TP+TN}{TP+TN+FP+FN},
\label{eq:evaluation-accuracy} \\
\mathrm{Precision} &= \frac{TP}{TP+FP},
\label{eq:evaluation-precision} \\
\mathrm{Recall} &= \frac{TP}{TP+FN},
\label{eq:evaluation-recall} \\
F_1 &= 2 \times
\frac{\mathrm{Precision}\times\mathrm{Recall}}
{\mathrm{Precision}+\mathrm{Recall}}.
\label{eq:evaluation-f1}
\end{align}

Accuracy represented overall predictive correctness, whereas precision measured the
proportion of predicted insecure files that belonged to the positive class. Recall
measured the proportion of labelled insecure files identified by the model. Recall was
particularly relevant to the gate because a false-negative result reduced the learned
risk contribution assigned to insecure code. The F1-score summarised the balance between
precision and recall, while ROC--AUC measured the ability of the model to rank positive
samples above negative samples across probability thresholds.

A scanner-only rule provided the baseline for comparison with the logistic-regression
classifier. The baseline classified a file as insecure when at least one configured
security finding was reported. This comparison determined whether learning a relationship
among the scanner features provided discriminatory value beyond treating the presence of
any finding as a positive classification.

\subsubsection{Unified Score and Decision Evaluation}
\label{subsubsec:evaluation-gate-decision}

The complete risk-scoring mechanism was evaluated using controlled cases containing
different combinations of deduplicated static findings, estimated infrastructure cost,
live compliance state, and machine-learning probability. For each case, the expected
component contributions and total score were calculated independently using
Equation~\ref{eq:score}. The independently calculated values were then compared with the
values recorded in the gate report.

Cases were constructed below, at, and above the decision boundaries defined in
Equation~\ref{eq:decision}. This boundary-oriented design exercised the \emph{pass},
\emph{review}, and \emph{reject} outcomes and examined whether the implemented
comparisons treated threshold values consistently. Additional cases contained
equivalent findings produced by more than one scanner to determine whether the
deduplication procedure prevented repeated counting of one underlying issue.

The gate evaluation considered three properties:

\begin{itemize}
    \item \textbf{Aggregation correctness}: the reported score equalled the sum of the
    deduplicated severity, cost, compliance, and machine-learning contributions.
    \item \textbf{Decision consistency}: the reported decision corresponded to the
    score interval defined by the configured thresholds.
    \item \textbf{Interpretability}: the report identified the contributing findings,
    component values, scanner execution states, total score, and resulting decision.
\end{itemize}

The classifier evaluation and the controlled gate cases addressed different aspects of
RQ3. The classifier measures assessed the discriminatory performance of the learned
code-risk component. The gate cases assessed whether that component and the other risk
signals were combined correctly and transparently to govern deployment.